\documentclass{article}
\input{../../../lib.sty}

\title{\texttt{fn group\_by\_output\_domain}}
\author{\MichaelShoemateAuthor}
\begin{document}
\maketitle

\contrib

\section{Hoare Triple}
\subsection*{Precondition}
\subsubsection*{Compiler-verified}
Types matching pseudocode.

\subsubsection*{Caller-verified}
The execution keys are stable row-by-row expressions with canonical representatives.
For every original key there is a corresponding execution key that preserves its values
and output name on the input domain, so the original margins describe the same partition.
The group-by uses ordinary options and no apply or having clause.

\subsection*{Pseudocode}
\lstinputlisting[language=Python,firstline=2,escapechar=|]{./pseudocode/group_by_output_domain.py}

\subsection*{Postcondition}
\begin{theorem}
If the constructor succeeds, ordinary group-by on the execution keys with the
given aggregations produces a member of the returned domain or raises a
data-independent runtime exception. Each aggregation depends only on its
own group and is invariant in distribution under permutations of its input rows.
Retained margins describe the corresponding materialized output keys.
\end{theorem}

\begin{proof}
    The caller supplies validated row-by-row grouping expressions with canonical representatives, preserving the values and output names of the original keys on the input domain.
    By line \ref{line:infallible-aggs} the aggregates are infallible.
    The check also excludes positional inputs and data-dependent quantile
    parameters, ensuring that aggregates depend only on their respective groups.
    On line \ref{line:permutation-invariance}, the checker certifies that
    row-aligned inputs commute with permutations and that order-sensitive
    reducers consume an invariant scalar or a trusted uniformly shuffled input.
    In particular, the shuffle distribution does not depend on the original
    physical row order.

    By the postcondition of \texttt{DslPlan.simulate\_schema},
    \texttt{series\_domains} follows the expected schema of members in the output domain.
    The output column domains preserve the schema while allowing infallible black-box aggregation.

    Only margins whose original grouping expressions are a subset of the grouping keys are preserved.
    On line \ref{line:materialized-margin}, each preserved key becomes \texttt{col(output\_name)}.
    Polars evaluates a grouping expression on the input and stores its value in the output column with that name.
    The translated margin therefore partitions output rows by the same labels as the original margin partitions input rows.
    Each input row contributes to at most one output row,
    so the upper bounds on group lengths and the number of groups remain valid.
    Re-evaluating the original expression after aggregation is not generally equivalent
    and must not inherit the original margin.
    Among these margins, invariants are discarded.

    It has been shown that for every element $x$ in \texttt{domain}, the group-by result on $x$ is in the returned domain or raises a data-independent runtime exception.
\end{proof}
\end{document}
